% !PRINTED 253
\weberpage{253}
%<<CONT

% [corrected 2026-08-07, 10x zoom strip 1/10] Two defects here, both found by the section-series
% check rather than by either audit lane.
%
% 1. A bare "\S.\,79." stood at the head of this fragment. It is not a heading -- it is the left
%    half of the RUNNING HEAD, "§. 79. | Bezoutiante und Wurzelrealität. | 253" -- and running
%    heads are excluded page furniture. The real §. 79 heading, with its title, is on printed 252
%    (i0277). Transcribing furniture as content gave the volume two §. 79 headings, which is how
%    it was caught: a heading with no title, sitting directly on body text, is the tell.
%
% 2. This theorem is letterspaced on the page and was set plain. It is the Satz announced by the
%    colon that ends printed 252 -- the same page whose missing %CONT>> the colon check flagged.
\gesperrt{Sind $x_1, x_2 \ldots x_\nu$ von einander verschiedene Wurzeln der
Gleichung (1), so sind $y_1, y_2 \ldots y_\nu$ linear unabhängig.}

Denn angenommen, es existire zwischen $y_1, y_2 \ldots y_\nu$ eine in
Bezug auf $t$ identische lineare Relation
\[
\alpha_1 y_1 + \alpha_2 y_2 + \cdots + \alpha_\nu y_\nu = 0,
\]
deren Coëfficienten $\alpha$ von den $t$ unabhängig und nicht alle gleich
Null sind, so würde diese in die folgenden Gleichungen zerfallen
\[
\begin{aligned}
\alpha_1 f_0(x_1) + \alpha_2 f_0(x_2) + \cdots + \alpha_\nu f_0(x_\nu) &= 0, \\
\alpha_1 f_1(x_1) + \alpha_2 f_1(x_2) + \cdots + \alpha_\nu f_1(x_\nu) &= 0, \\
\cdot\ \cdot\ \cdot\ \cdot\ \cdot\ \cdot\ \cdot\ \cdot\ \cdot\ \cdot\ \cdot\ \cdot\ \cdot\ \cdot & \\
\alpha_1 f_{n-1}(x_1) + \alpha_2 f_{n-1}(x_2) + \cdots + \alpha_\nu f_{n-1}(x_\nu) &= 0.
\end{aligned}
\]
Da nun $\nu \leqq n$ ist,
so muss die Determinante der $\nu$ ersten
von diesen Gleichungen verschwinden, also
\[
\begin{vmatrix}
f_0(x_1), & f_1(x_1) & \ldots & f_{\nu-1}(x_1) \\
f_0(x_2), & f_1(x_2) & \ldots & f_{\nu-1}(x_2) \\
\cdot & \cdot & \cdot\ \cdot\ \cdot\ \cdot\ \cdot & \cdot \\
f_0(x_\nu), & f_1(x_\nu) & \ldots & f_{\nu-1}(x_\nu)
\end{vmatrix}
= 0,
\]
oder
\[
\begin{vmatrix}
1, & x_1 + a_1, & x_1^2 + a_1 x_1 + a_2 \ldots \\
1, & x_2 + a_1, & x_2^2 + a_1 x_2 + a_2 \ldots \\
\cdot & \cdot & \cdot\ \cdot\ \cdot\ \cdot\ \cdot\ \cdot\ \cdot\ \cdot \\
1, & x_\nu + a_1, & x_\nu^2 + a_1 x_\nu + a_2 \ldots
\end{vmatrix}
= 0.
\]
Diese Determinante lässt sich aber durch wiederholte Anwendung
des Satzes von der Addition der Colonnen \S.\,22, (VII)
auf die Form bringen
\[
\pm
\begin{vmatrix}
1, & x_1, & x_1^2 & \ldots & x_1^{\nu-1} \\
1, & x_2, & x_2^2 & \ldots & x_2^{\nu-1} \\
\cdot & \cdot & \cdot & \cdot\ \cdot\ \cdot\ \cdot & \cdot \\
1, & x_\nu, & x_\nu^2 & \ldots & x_\nu^{\nu-1}
\end{vmatrix}
=
\begin{aligned}
&(x_1 - x_2)(x_1 - x_3) \ldots (x_1 - x_\nu) \\
&\hphantom{(x_1 - x_2)}(x_2 - x_3) \ldots (x_2 - x_\nu) \\
&\hphantom{(x_1 - x_2)(x_2 - x_3)}\cdot\ \cdot\ \cdot\ \cdot\ \cdot \\
&\hphantom{(x_1 - x_2)(x_2 - x_3)\cdot\cdot}(x_{\nu-1} - x_\nu),
\end{aligned}
\]
und kann also nicht verschwinden, wenn, wie vorausgesetzt war,
die $x_1, x_2 \ldots x_\nu$ verschieden sind.

Wenn wir daher in der Summe (4) alle unter einander gleichen
Glieder zusammenfassen, \gesperrt{so bleiben so viele Quadrate
linear unabhängiger Functionen übrig, als die Gleichung (1)
verschiedene Wurzeln hat.}
